\begin{figure}[htbp]
  \centering
  \resizebox{\textwidth}{!}{%
  \begin{tikzpicture}[
    >=Latex,
    font=\small,
    node distance=10mm and 14mm,
    class/.style={
      draw,
      rounded corners=3pt,
      align=left,
      fill=gray!8,
      minimum width=42mm
    },
    support/.style={
      draw,
      rounded corners=3pt,
      align=left,
      fill=blue!6,
      minimum width=34mm
    },
    arrow/.style={-{Latex[length=2.3mm]}, thick}
  ]
    \node[class] (generator) {
      \textbf{Generator}\\
      create()\\
      evaluate()\\
      export()
    };

    \node[class, left=24mm of generator] (melody) {
      \textbf{Melody}\\
      title\\
      key\\
      events: Note[]
    };

    \node[support, below=of melody] (note) {
      \textbf{Note}\\
      pitch\\
      duration
    };

    \node[support, right=24mm of generator] (constraint) {
      \textbf{Constraint}\\
      check(candidate)
    };

    \node[support, above right=of constraint] (range) {
      \textbf{RangeConstraint}\\
      checks pitch range
    };

    \node[support, below right=of constraint] (cadence) {
      \textbf{CadenceConstraint}\\
      checks phrase ending
    };

    \draw[arrow] (generator) -- node[above] {creates} (melody);
    \draw[arrow] (note) -- node[left] {contained in} (melody);
    \draw[arrow] (constraint) -- node[above] {used by} (generator);
    \draw[arrow] (range) -- (constraint);
    \draw[arrow] (cadence) -- (constraint);
  \end{tikzpicture}%
  }
  \caption{A simple object-oriented example showing two important ideas. Composition is used when a melody consists of note objects, while polymorphism is used when different constraint classes can be treated through one shared interface.}
  \label{fig:oop_iter3_model}
\end{figure}
